\documentclass[11pt,a4paper]{article}
\usepackage[T1]{fontenc}
\usepackage[utf8]{inputenc}
\usepackage{lmodern}
\usepackage[margin=27mm]{geometry}
\usepackage{amsmath,amssymb,amsthm}
\usepackage{array,booktabs,tabularx}
\usepackage{enumitem}
\usepackage[hidelinks]{hyperref}
\usepackage{microtype}
\setlength{\parindent}{0pt}
\setlength{\parskip}{5pt plus 1pt}
\setlist[itemize]{leftmargin=1.45em,itemsep=2pt,topsep=3pt}
\setlist[enumerate]{leftmargin=1.7em,itemsep=2pt,topsep=3pt}
\newcommand{\Ctx}{\operatorname{Ctx}}
\newcommand{\Cnt}{\operatorname{Cnt}}
\newcommand{\Cat}{\mathbf{Cat}}
\newcommand{\id}{\mathrm{id}}
\newcommand{\FinStoch}{\mathbf{FinStoch}}
\newcommand{\HistProb}{\mathbf{HistProb}}
\title{Generalizing contextads\\[4pt]\large Provisional research contract for dependent mixed composition}
\author{}
\date{Draft for discussion\quad 1 November 2026}

\begin{document}
\maketitle
\vspace{-1.5em}

\textbf{Status.} Drafting has been authorized. The research agreement, execution
of the proposed workstream and its final theorem hypotheses have not yet been
approved.

\textbf{Decision requested.} Consider a bounded first workstream on dependent mixed
context and content composition, with a limited recognition or necessity
component. Begin by auditing the strongest relevant results and typing the
candidate construction. Only then agree its theorem hypotheses. This draft
sets desired mathematical content and acceptance tests; it does not certify
novelty, promise a universal theorem, or freeze the proposed assumptions.

\section{Programme and first workstream}

The programme is to \emph{generalize contextads}, through theoretical results
and multiple applications. It includes three different kinds of progress:
instantiating existing contextads, proving new construction or comparison
results, and extending the ambient machinery or coherence requirements.
The research must say which kind a result achieves. A richer model within the
existing definition is valuable but is not automatically a new definition.

The first workstream asks when a contextad and a contentad in a declared
common setting admit coherent mixed composition: a computation requires
context, produces content, and its realized result can change the base over
which the next requirement is indexed. The principal target is one abstract
family theorem, not a composition formula for one finite system. A bounded
necessity or recognition statement should explain what the selected
representation requires; reconstructing every contextad from arbitrary double
categories is outside this contract.

Transactions by robots with discovered semantics and partial undo are one
demanding application, not the definition of the programme or its first theorem.
Adaptive probabilistic programming and typed protocols with dynamically
allocated capabilities provide independent tests. Differentiable programming
is another candidate application. Noncartesian ambient constructions, an
internal language, higher order, differentiation and feedback remain possible
later directions, not automatically contracted features.

\textbf{Desired content held fixed for this draft.} The work must address
actual dependence of later contexts on produced results; recover the existing
one-sided and fixed-functor controls; and supply a usable compatibility result
with a sharp boundary. These goals cannot be replaced silently by a theorem
whose only hypothesis is the desired mixed law, by a single worked example,
or by an unrelated weakening that is easier to prove.

\textbf{Decisions still provisional.} The ambient category or 2-category,
display and codisplay classes, admissible morphisms, representation of mixed
arrows, equality level, and any factorization or exactness assumptions remain
to be selected. An indexed presentation in \(\Cat\), tested first on finite
examples, is a starting candidate, not a requirement that all applications be
finite or probabilistic. The audit may recommend a different starting class,
with reasons and an explicit scope decision.

\section{What the present evidence already establishes}

Capucci and Myers already cover input-dependent contexts, noninvertible
context-extension structure maps, partial-map and state examples, and
2-algebraic structure lifting under stated hypotheses. Their Section 7 permits
coherent transport without universal cartesian lifting, but retains invertible
structural comparisons. Section 8 explicitly proposes mixed contextad and
contentad laws and factorization-based recognition; these are not new proposals
of this project \cite{CM}. An equivalent indexed presentation alone is not a
generalization.

The companion finite report supplies a useful positive control: dependent
families of finite probability spaces form an ordinary contextad. Evaluation
\(E:\HistProb\to\FinStoch\) forgets sample histories. Its canonical section
\(J\) is locally fully faithful, with local adjunctions \(E_{A,B}\dashv J_{A,B}\)
and \(EJ=\id\), but its global composition comparison generally forgets
intermediate samples and is not invertible \cite{Finite}. Thus local reflection,
law equality and history-preserving equivalence must not be conflated.

That report also proves a narrow impossibility for deterministic writer grades
paired with extensional normalized probability laws while preserving the
pure-writer and grade-zero stochastic inclusions. Random grades or retained sample presentations repair
that example. Neither it nor the failure of one mixture-transport formula is
a universal obstruction. These results are starting controls, not by
themselves completion of the new workstream.

\section{Seven substantive obligations}

\subsection{Audit the strongest baselines and agree a typed charter}

Before fixing theorem obligations, inspect the closest primary definitions,
results and proofs, not only abstracts or statements of future work. The audit
must compare the proposed result with contextads and contentads, generalized
Section 7 transport, and applicable pseudodistributive-law or lifting theory.
It must also check dependent sum or factorization constructions and any
polynomial or parametric-right-adjoint transposition invoked by the proposed
representation. Walker's presentations are a relevant entry point, not a
result already specialized to this problem \cite{Walker}.

Record each material overlap as one of: already proved and applicable;
applicable only after additional hypotheses; stated or conjectured without a
verified proof; or still unresolved in this audit. In particular, the
polynomial transposition statement in Appendix A of the 2024 paper is
explicitly left unproved there. The 2025 slides present a more general PRA
statement as a theorem, but the preliminary audit did not verify a complete
proof \cite{CM,Slides}. Neither blanket dismissal nor unqualified import of
that result is warranted.

The resulting short charter must specify the common base; the two input
structures and their maps; mixed objects, arrows and comparison cells; the
permitted continuation class; and the intended equality or equivalence level.
It must identify every proposed pullback, pushforward, dependent sum, strength,
choice of presentation or factorization used to make composition type-check.
Every substantial hypothesis gets a role: needed for typing, used in a proof,
imported from an existing theorem, or conjecturally relaxable. State whether
coherence is derived or supplied as data.

\textbf{Acceptance.} A reader can see the exact increment over the strongest
verified baseline and the unresolved assumptions. Review and agree this charter
before treating any theorem statement as fixed. Finding that the proposed
result is already known is a valid audit outcome, but triggers an explicit
choice of a distinct target; it is not disguised as new research.

\subsection{Construct dependent mixed composition and prove its coherence}

For the agreed class, give the complete typed construction, including identities,
composition and the morphisms or cells needed for comparison. Locate the
produced result and the next context in their actual fibres. Show how the
interaction transports or combines them; an untyped expression with a context
on the left and content on the right does not suffice. State the ambient
structure actually obtained, such as a category, bicategory or pseudocategory.

Prove the unit and associativity laws at that level, with the required
naturality and coherence. Where cartesian or opcartesian preservation is
claimed, prove it rather than inferring it from an object-level formula.
Prove functorial comparison under the declared maps of examples, at the
stated strict, pseudo or lax level. Do not imply a higher-dimensional universal
property beyond the one proved.

\textbf{Acceptance.} The result covers a parameterized class and can be
instantiated without inventing its composition anew. All transports and
structural comparisons have explicit sources, targets and directions. A
construction retaining full histories is admissible, but must state what it
adds and what its comparison forgets; merely enriching the base does not
establish the sought compatibility criterion.

\subsection{Give a usable compatibility criterion and bounded necessity}

The substantive result must go beyond ``given a suitable mixed distributive
law, composition is associative.'' That is a formal baseline. For the chosen
class, derive a nontrivial interaction or its coherence from independently
checkable structural conditions, such as a specified universal property or
stable factorization. Separate these conditions from any higher coherence
which must still be supplied.

Add a limited recognition or necessity statement within an explicitly stated
normal-form or preservation class. For example, prove that a proposed
representation preserving specified pure arrows and observations forces a
particular interaction equation or factorization, and explain its relation to
the sufficient criterion. Necessity outside that class is not claimed. A
complete intrinsic characterization of all mixed structures, or globally
weakest hypotheses, is not required.

\textbf{Anti-tautology test.} Apply the criterion to independently described
data and show how it can be checked without already possessing the desired
mixed composition. Exhibit which genuine coherence or compatibility fact the
proof derives. If ``good factorization'' simply packages the target structure,
label the result a reformulation and do not count it as this obligation.
If generic pseudodistributive-law theory proves everything claimed, identify
the extra specialization, recognition or boundary result that remains.

\textbf{Acceptance.} There is a noncircular entry test, a proved consequence,
and at least one nontrivial necessary condition or reconstruction fact relative
to the frozen class. No unrestricted recognition theorem is promised.

\subsection{Recover the conservative controls at declared equality levels}

The construction must not change the meanings of its familiar limits. Prove
the following comparisons, including their action on pure or tight arrows and
on the relevant cells.

\begin{center}
\begin{tabularx}{\linewidth}{@{}p{0.23\linewidth}X@{}}
\toprule
Control & Required comparison \\
\midrule
Context only & With content trivialized as specified, recover \(\Ctx\) by a
structure-preserving equivalence over the base. Claim an isomorphism or literal
equality only when established for the chosen presentation. \\
Content only & Dually, with context trivialized, recover \(\Cnt\) with the same
explicit standard of comparison. \\
Fixed functors & For an ordinary comonad \(G\), monad \(T\), and compatible
mixed law \(\lambda:GT\Rightarrow TG\), recover the ordinary mixed arrow
category and its composition, with equality of the displayed arrow formulas. \\
\bottomrule
\end{tabularx}
\end{center}

For the fixed-functor control, if \(f:GA\to TB\) and \(g:GB\to TC\), the
required formula is
\[
 g\star f
 =\mu_C\circ T g\circ\lambda_B\circ G f\circ\delta_A,
 \qquad
 1_A=\eta_A\circ\varepsilon_A.
\]
Here \((G,\varepsilon,\delta)\) and \((T,\eta,\mu)\) have their usual
comonad and monad structures. The compatibility axioms on \(\lambda\) must be
stated. A pseudo analogue, if used, has its own coherence requirements and is
not inferred from this ordinary-category calculation.

Any passage to extensional laws, connected components or another quotient
must be separately identified. State which distinctions and pure arrows it
preserves, and whether the comparison is faithful, an equivalence, reflective,
or merely lax. An isomorphism of outputs is not an isomorphism of retained
histories, and a noninvertible comparison cannot be renamed an equivalence.

\textbf{Acceptance.} The limiting constructions are conservative at exactly
the claimed level, and the report explains when an ordinary contextad
presentation suffices and what additional distinctions the mixed presentation
retains. Compare alternative encodings relative to declared preservation of pure arrows,
context and content distinctions, or continuation interfaces. Show what a
fixed-control encoding must enlarge or forget to represent the dependent case.
No absolute impossibility of one-sided encoding or of flattening into a larger
state space is required.

\subsection{Test contrasting applications with genuine dependence}

Instantiate the abstract result in at least two contrasting application
families. At least one example must be genuinely dependent: an intermediate
result changes the next context's fibre, admissible type or available
structure, and this dependence appears in the composition proof rather than
only in the names of objects. At least one must be non-robotic and require no
recovery. Explain the distinct requirement each example tests.

The provisional working pair is adaptive probabilistic computation and a
deterministic typed protocol with dynamic capability allocation. The first
can distinguish shared latent information, fresh samples and output laws;
the second can test result-indexed legality and resource creation without
probability. Selection and exact models are to be agreed at the audit gate;
this is not a promise to solve the general semantics of either application.
Ordinary state, indexed or graded effects and existing protocol semantics
remain comparison baselines.

The robot example may provide an additional stress test. If included, separate
latent execution witnesses, observable receipts and allowed continuations.
No proof witness grants the controller an observation, permission, reset or
conditioning operation. State the recovery domain and residual correctness
relation rather than equating compensation with inversion. Discovered semantics
may invalidate an earlier model; they do not automatically give a lawful
refinement. A general learning or partial-undo calculus is not required here.

\textbf{Acceptance.} Both examples verify the criterion, exercise the theorem
and identify their operational or representational limits. At least one
nonidentity change of instance or representation must exercise the naturality
result through its dependence or interaction, unless the audit gate expressly
agrees another substantive witness. Reusing the
finite probability-family baseline is welcome as a control, but that baseline
alone does not discharge the mixed or multi-application obligations.

\subsection{Prove a sharp obstruction and a useful positive branch}

Identify a plausible overextension or omitted hypothesis of the proposed
construction. Prove a counterexample or no-go result with the class of
representations, allowed operations, preserved embeddings and equality level
stated explicitly. Show why the example satisfies the retained assumptions
and violates the claimed conclusion. Distinguish failure of existence,
failure of preservation and loss of presentation information.

Pair that obstruction with a useful positive branch: strengthen a hypothesis,
change a representation, or retain necessary information, then prove precisely
what works and what the repair sacrifices. The repair must connect to the
contracted family theorem and still admit genuine dependent behavior. Repeating
the known finite writer obstruction and its standard random-grade repair is
background, not the new boundary result.

\textbf{Acceptance.} The obstruction tests the theorem's real boundary rather
than an unrelated toy encoding. The repair is proved, not merely proposed, and
its scope cost is explicit. A negative result ruling out the proposed main
branch is useful research but not automatic completion: if no adequate
positive branch remains, present the obstruction and a proposed rescope for
decision instead of silently substituting a weaker task.

\subsection{Deliver a proof dossier with honest evidence status}

Deliver a self-contained mathematical report containing the agreed definitions,
main family theorem and proof, compatibility and limited necessity result,
conservative comparisons, application instances, and obstruction with repair.
Integrate a compact assumption and source ledger: for each important
hypothesis, cite its use; for each imported result, give its source and verified
proof status; for each claimed advance, state the closest overlap. Include a
short account of changed hypotheses and unresolved questions.

Mark assertions as proved, checked only in specified finite cases,
conjectural, or unverified. Computation can find counterexamples and check
finite instances; it does not replace a parametric proof. If a finite exhaustive
check itself proves a bounded claim, state the exact bound and give a
reproducible procedure. Proof-assistant formalization is optional and must
report precisely what is formalized; no Lean implementation or TeX compilation
is a completion requirement of this contract.

\textbf{Acceptance.} Another researcher can reconstruct the result, check its
assumptions and instantiate its examples without relying on hidden code,
unstated coherence or unverified literature. A result already known, an
incomplete proof and a new proved theorem are reported differently. This
contract promises a disciplined investigation, not a novelty conclusion in
advance.

\section{Review gates and stopping condition}

\textbf{First gate.} The proposed next step, subject to authorization, is scoping
and the strongest-source audit, not adoption of every candidate axiom. The audit
ends with the typed charter and a recommended theorem package. The user reviews
that package before its hypotheses and success criteria are frozen. Exploratory
calculations within an authorized scoping stage may inform the choice, but do
not commit it.

\textbf{Change control.} A changed ambient class, weaker preservation claim,
coarser equality, restricted continuation language or less dependent example
must be recorded with its mathematical cause and consequence. Present the old
and proposed statements together and seek agreement for a material reduction
or replacement. A sharper assumption discovered in proof is not concealed in
a definition. A literature result that subsumes the target prompts an explicit
reassessment rather than a novelty claim based on renaming.

\textbf{Completion.} The bounded workstream is complete when the agreed seven
obligations have been discharged and the dossier clearly states what was
proved, what was inherited, and what remains open. It stops at that point;
there is no automatic move to every Section 8 frontier or every candidate
application. If a central obligation cannot be met, deliver the exact blocking
result or gap and the next decision needed. That is an honest paused or
rescoped outcome, not a claim that the original theorem package was fulfilled.

\begin{thebibliography}{9}
\bibitem{CM}
Matteo Capucci and David Jaz Myers.
\emph{Contextads as Wreaths; Kleisli, Para, and Span Constructions as Wreath
Products}. arXiv:2410.21889v1, 2024.
Definitions 3.15 and 3.26; Sections 5--8; Appendix A.
\url{https://arxiv.org/abs/2410.21889}.

\bibitem{Walker}
Charles Walker.
\emph{Presentations of pseudodistributive laws}.
arXiv:2102.12468v4.
A literature-audit target; no specialized mixed-contextad theorem is imported
here.
\url{https://arxiv.org/abs/2102.12468v4}.

\bibitem{Slides}
Matteo Capucci.
\emph{Contextads abridged}. PSSL110 slides, May 2025.
PRA construction and stated Kleisli--Ctx result; a complete proof was not
verified in the preliminary source audit.
\url{https://matteocapucci.eu/assets/2025/05/pssl-110.pdf}.

\bibitem{Finite}
\emph{Generalizing contextads: A finite stochastic baseline and one recovery
application}. Project report, scope-corrected 1 November 2026.
Sections 4--8 contain the finite writer obstruction, probability-family
contextad and history-to-law comparison used as controls here.
\end{thebibliography}

\end{document}
